\section{Conclusion}\label{sec:conclusion}

Logs that record edits but not reads cannot say how a mistake spread through a swarm: in a real incident, even a perfect investigator could name the source of at most \bestTop{} of copied behaviours, and the obvious detector mostly produced false calls. We characterised what can and cannot be recovered, called copying against base rates, and showed that a gateway that serves every read with its own load-bearing token turns a copied link into a pointer to the exact copy it came from. The benefit is not uniform: it follows how often a model copies links instead of re-deriving access, and it becomes uniform only when the token is made the only way in. Tracing a planted bad tip shrank what had to be re-checked from most of the outputs to under a tenth, unless the tip was there from the start. We hope that the audit, the gateway design and the benchmark with its hidden read log are useful to people who run swarms, and that the list of what went wrong saves them time. The first thing to try is cheap: run the audit on the logs you already keep.
